\chapter{Conditional expectation and linear projection}
\label{cap:ec}

\section{What conditional expectation is for}
Almost all of econometrics tries to answer: \emph{given what I know about $x$, what is the typical value of $y$?} The mathematical answer is the conditional expectation $m(x)=\E(y\mid x)$.

\begin{intuicao}
Imagine splitting people by years of schooling ($x$) and computing the average wage of each group ($y$). The list of those group averages is $m(x)$. It says nothing about cause. It only says what the average wage is for people with that schooling.
\end{intuicao}

\begin{resultado}[Conditional expectation is the best predictor]
For any function $g(x)$,
\[
\E\big[(y-g(x))^2\big] \;=\; \E\big[(y-m(x))^2\big] \;+\; \E\big[(m(x)-g(x))^2\big] \;\ge\; \E\big[(y-m(x))^2\big].
\]
Hence $m(x)$ minimizes the mean squared error among all functions of $x$.
\end{resultado}
\emph{Why?} Add and subtract $m(x)$ inside the square. The cross term has zero expectation because $\E(y-m(x)\mid x)=0$.

\section{The error form}
Define $u = y - m(x)$. Then $y=m(x)+u$ with $\E(u\mid x)=0$. Two consequences used all the time:
\begin{itemize}[leftmargin=1.4em]
\item $\E(u)=0$ (by the law of iterated expectations, $\E(u)=\E[\E(u\mid x)]$);
\item $\Cov(u,h(x))=0$ for any function $h$ (the error has no relation with $x$, neither linear nor nonlinear).
\end{itemize}

\section{Partial effects, elasticities and semi-elasticities}
If $m(x_1,x_2)$ is the conditional expectation, the partial effect of $x_1$ is $\partial m/\partial x_1$. With variables in logs:
\begin{center}
\small
\begin{tabular}{lll}
\toprule
\textbf{Model} & \textbf{Coefficient $\beta_1$ measures} & \textbf{Name} \\
\midrule
$y=\beta_0+\beta_1x_1$ & change in $y$ per unit of $x_1$ & marginal effect \\
$\log y=\beta_0+\beta_1 x_1$ & $\approx 100\beta_1\%$ change in $y$ per unit of $x_1$ & semi-elasticity \\
$\log y=\beta_0+\beta_1\log x_1$ & $\%$ change in $y$ per $1\%$ in $x_1$ & elasticity \\
\bottomrule
\end{tabular}
\end{center}
If the model has an interaction or a quadratic term, the effect of $x_1$ changes with the values of the variables. In that case report the \textbf{average partial effect} (APE): the average of the partial effect over the distribution of $x$.

\section{Linear projection: what OLS really estimates}
$m(x)$ is not always linear. Even so, there is always a best linear approximation in the mean squared error sense:
\[
L(y\mid 1,x)=\beta_0+x\beta,\qquad \beta=\big[\E(x'x)\big]^{-1}\E(x'y)\ \ (\text{with } x \text{ including the constant}).
\]
\begin{intuicao}
The linear projection is the line (or hyperplane) that best \emph{summarizes} the cloud of points, even if the cloud is curved. OLS estimates exactly this line. Only when $m(x)$ really is linear does this line coincide with the conditional expectation.
\end{intuicao}

\begin{atencao}[Pitfall: ``linear'' is not ``causal'']
The projection coefficient is always well defined. That does not make it a causal effect. For $\beta$ to have a causal reading, you need an extra argument about why the other factors that affect $y$ do not move together with $x$ (chapters \ref{cap:mqo} to \ref{cap:vi}).
\end{atencao}

\begin{exercicio}
Let $x\sim N(0,1)$ and $y=x^2$ (no noise). What is the slope of the linear projection of $y$ on $x$? What does this say about dependence between $y$ and $x$?
\end{exercicio}
\sol{$\Cov(x,x^2)=\E(x^3)-\E(x)\E(x^2)=0$, because odd moments of a centered normal are zero. So the slope is zero and the projection is the constant $\E(y)=1$. Yet $y$ is an exact function of $x$. A zero correlation (or slope) does \textbf{not} mean independence. This is the shortest example of why the line can hide a strong relationship.}

\begin{exercicio}
Show that if $\E(u\mid x)=0$, then $\Cov(u,x^2)=0$.
\end{exercicio}
\sol{$\Cov(u,x^2)=\E(ux^2)-\E(u)\E(x^2)$. By the law of iterated expectations, $\E(ux^2)=\E[x^2\E(u\mid x)]=0$, and $\E(u)=0$. Therefore the covariance is zero.}
